\hypertarget{appendix-ce20506c8532}{}
\section*{Signed magnitude}
\section{Signed magnitude}\label{signed-magnitude}

\textbf{Curriculum depth:} Introduced · \textbf{First seen:} C002 · \textbf{Latest development:} C002

\subsection{Definition and mechanism}\label{definition-and-mechanism}

The high bit selects sign while remaining bits encode magnitude; this is intuitive but has positive and negative zero and needs sign-aware arithmetic.

The high bit holds sign and the remaining bits hold magnitude. Eight-bit +5 is \texttt{00000101}, −5 is \texttt{10000101}, and \texttt{00000000} and \texttt{10000000} both denote zero. Mixed signs require magnitude comparison and subtraction rather than ordinary raw-bit addition.

\subsection{Prerequisites}\label{prerequisites}

\begin{itemize}
\tightlist
\item
  \hyperref[appendix-83c3f28f7bda]{Positional notation} --- A base is the number of available digits.
\item
  \hyperref[appendix-280445746493]{Bit significance} --- LSB means least significant bit, normally position 0 with weight 1.
\end{itemize}

\subsection{Used by}\label{used-by}

\begin{itemize}
\tightlist
\item
  \hyperref[appendix-03c678cd40b3]{One\textquotesingle s complement} --- uses signed magnitude as part of its representation or reasoning.
\end{itemize}

\subsection{Recall and next depth}\label{recall-and-next-depth}

Explain signed magnitude without the label, work the example above, then state the assumption whose removal would change the result.

\textbf{Pending:} floating-point sign fields and non-integer signed encodings.

\subsection{Source and navigation}\label{source-and-navigation}

\href{Continuous\%20Notes/Computing\%20Foundations\%20-\%20From\%20Bits\%20to\%20C++\%20and\%20Systems\%20APIs\#Arithmetic\%20becomes\%20a\%20policy}{C002 teaching chapter} · \hyperref[appendix-fd090d982a15]{Part definitions and mechanisms} · \href{}{Dependency Map (02-Foundations-II, appendix-4ad47bdd269a)}
